% Auto-split to keep each TeX source < 800 lines.

\begin{corollary}[最大纤维、组合上界与饱和判别]\label{cor:fold-bin-saturation}
对任意 \(m\ge 1\)，最大纤维由 \(w=0^m\) 取得，并满足组合上界
\begin{equation}\label{eq:fold-bin-upper}
d^{\mathrm{bin}}_{\max}(m):=d^{\mathrm{bin}}_m(0^m)\le |X_{K(m)-m}|=F_{K(m)-m+2}.
\end{equation}
记 \(K=K(m)\)、\(L=K-m\)，并定义在相邻约束下的\emph{最大尾和}
\[
S_{\max}(m):=\max\Bigl\{\sum_{k=m+1}^{K} c_k F_{k+1}:\ (c_{m+1},\dots,c_K)\in\{0,1\}^{L},\ c_k c_{k+1}=0\Bigr\}.
\]
则对 \(w=0^m\)，上界 \eqref{eq:fold-bin-upper} 饱和当且仅当阈值约束对所有合法尾部都是冗余的，即
\begin{equation}\label{eq:fold-bin-sat-iff}
d^{\mathrm{bin}}_{\max}(m)=F_{K(m)-m+2}\quad\Longleftrightarrow\quad 2^m-1\ge S_{\max}(m).
\end{equation}
并且 \(S_{\max}(m)\) 具有完全显式的闭式（由“1010\(\cdots\)”极大尾部取得）：
\begin{equation}\label{eq:fold-bin-Smax-closed}
S_{\max}(m)=
\begin{cases}
F_{K(m)+2}-F_{m+1}, & L=K(m)-m\ \text{为奇数},\\
F_{K(m)+2}-F_{m+2}, & L=K(m)-m\ \text{为偶数}.
\end{cases}
\end{equation}
等价地，饱和判别可改写为
\[
F_{K(m)+2}-2^m\le
\begin{cases}
F_{m+1}-1,& L\ \text{奇},\\
F_{m+2}-1,& L\ \text{偶}.
\end{cases}
\]
\end{corollary}

\begin{theorem}[最大纤维指数与 \texorpdfstring{$\log(2/\varphi)$}{log(2/phi)}（bin-fold）]\label{thm:fold-bin-max-fiber-exponent}
令 \(d^{\mathrm{bin}}_{\max}(m):=\max_{w\in X_m} d^{\mathrm{bin}}_m(w)\)。
则极限存在并满足
\[
\lim_{m\to\infty}\frac{1}{m}\log d^{\mathrm{bin}}_{\max}(m)=\log\frac{2}{\varphi}.
\]
特别地，该指数常数与正文命题 \ref{prop:beta-quantization-step} 中的不可约步长 \(\Delta_{\min}=\log(2/\varphi)\) 同一。
\end{theorem}
\begin{proof}
由满射 \(\mathrm{Fold}^{\mathrm{bin}}_m:\{0,\dots,2^m-1\}\twoheadrightarrow X_m\) 的计数恒等式
\(
\sum_{w\in X_m} d^{\mathrm{bin}}_m(w)=2^m
\)
得
\(
d^{\mathrm{bin}}_{\max}(m)\ge 2^m/|X_m|
=2^m/F_{m+2}
\)。
由 Binet 渐近 \(F_{m+2}=\Theta(\varphi^{m})\) 得
\[
\liminf_{m\to\infty}\frac{1}{m}\log d^{\mathrm{bin}}_{\max}(m)\ \ge\ \log\frac{2}{\varphi}.
\]
\par
另一方面，由推论 \ref{cor:fold-bin-saturation} 的上界
\(
d^{\mathrm{bin}}_{\max}(m)\le F_{K(m)-m+2}
\)。
由定义 \eqref{eq:K-of-m} 有
\(
F_{K(m)+1}\le 2^m-1 < F_{K(m)+2}
\)，
结合 \(\log F_n=n\log\varphi+O(1)\) 得
\[
\bigl(K(m)+1\bigr)\log\varphi+O(1)\ \le\ m\log 2\ \le\ \bigl(K(m)+2\bigr)\log\varphi+O(1),
\]
从而 \(K(m)=m\log_{\varphi}2+O(1)\)，故
\[
\log F_{K(m)-m+2}=\bigl(K(m)-m\bigr)\log\varphi+O(1)
=m(\log 2-\log\varphi)+O(1).
\]
两边同除以 \(m\) 并取上极限，得
\(
\limsup_{m\to\infty}\frac{1}{m}\log d^{\mathrm{bin}}_{\max}(m)\le \log(2/\varphi)
\)。
合并上下界即得极限同一性。证毕。
\end{proof}

\begin{remark}[指数级 Diophantine 近似与“仅有限次饱和”]\label{rem:fold-bin-finite}
由 \eqref{eq:fold-bin-Smax-closed} 可见，饱和条件强迫 \(2^m\) 与某个 Fibonacci 数 \(F_{K(m)+2}\) \emph{异常接近}。
用 Binet 公式 \(F_n=\varphi^n/\sqrt5+O(\varphi^{-n})\) 将其改写，可导出一个指数级小的线性对数形式约束：
\[
\Bigl|(K(m)+2)\log\varphi - m\log 2 - \log\sqrt5\Bigr|
\ \lesssim\ \Bigl(\frac{\varphi}{2}\Bigr)^m.
\]
由于 \(\varphi\) 为代数数而 \(2\) 为有理数，上式左侧是非零的对数线性形式；而 Baker 型下界（例如 Matveev 的显式形式）给出其\emph{至多多项式级}的小下界（见 \cite{Matveev2000LinearFormsLogs,BakerWustholz2007LogarithmicForms}），从而与右侧的指数级小量矛盾，推出：
\[
\text{上界饱和只能发生有限次。}
\]
本结论在本文的可复现管线中可直接数值审计（见表 \ref{tab:fold_binary_saturation_points}）。
\end{remark}

\begin{proposition}[计算 \(\boldsymbol{d^{\mathrm{bin}}_m(w)}\) 的 \(O(m)\) Fibonacci-digit DP]\label{prop:fold-bin-digit-dp}
给定 \(m\ge 1\) 与 \(w\in X_m\)。令 \(K=K(m)\) 并记 \(V=V_m(w)\)，再令 \(\mathrm{slack}:=2^m-1-V\)。
将 \(\mathrm{slack}\) 写为 Zeckendorf 位串 \(Z(\mathrm{slack})=(b_1,\dots,b_K)\)（合法，无相邻 \(11\)）。
则 \(d^{\mathrm{bin}}_m(w)\) 可由如下 digit DP 在 \(O(K-m)=O(m)\) 时间内精确计算：
\[
\mathrm{DP}(k,\mathrm{prev},\mathrm{tight})
\quad(k=K,K-1,\dots,m+1),
\]
其中 \(\mathrm{prev}\in\{0,1\}\) 记录上一个（更高位）是否取 \(1\)，\(\mathrm{tight}\in\{0,1\}\) 表示当前前缀是否仍与 \((b_k,\dots,b_K)\) 相同；
转移时允许选择 \(c_k\in\{0,1\}\)，满足
\[
c_k c_{k+1}=0,\qquad
\mathrm{tight}=1\Rightarrow c_k\le b_k,\qquad
w_m=1\Rightarrow c_{m+1}=0,
\]
并按是否保持等号更新 \(\mathrm{tight}\)。该 DP 的输出即为满足 \(\sum_{k=m+1}^K c_k F_{k+1}\le \mathrm{slack}\) 的尾部数目，从而等于 \(d^{\mathrm{bin}}_m(w)\)。
\end{proposition}

\input{sections/generated/tab_fold_binary_saturation_points}

